% =============================================================================
% SECTION 5: ALP BIREFRINGENCE — THE SURVIVING PREDICTION
% =============================================================================
% Target: ~3.5 pages PRD format
% Framing: payoff of the closure assessment, not standalone ALP paper.
% =============================================================================

\section{Spectator ALP Birefringence}
\label{sec:alp}

The structural closure of Sec.~\ref{sec:closure} eliminates all direct routes
from the bounce to late-time observables.  One indirect testable handle
remains: the ECH framework's parity-odd sector motivates an effective
axion-like particle at the Planck scale, whose coupling to photons produces
cosmic birefringence.  We stress at the outset that this ALP is not uniquely
derived from ECH---any Planck-scale pseudoscalar with Standard Model anomaly
coupling yields the same prediction~\cite{Fujita2021,Nakagawa2025}.  The ECH
framework provides one possible UV completion; our contribution is the
identification of this prediction as the sole survivor of a comprehensive
assessment, and its confrontation with data through MCMC analysis.


% ---------------------------------------------------------------------------
\subsection{The Effective Spectator ALP}
\label{sec:alp_model}

The ECH parity-odd sector generates, at one loop, a finite parity-violating
coefficient of natural order $\mathcal{O}(\alpha_\mathrm{em}/M_\mathrm{Pl})$
\cite{Alexander2009,Freidel2005}.  We parametrize this sector as an effective
pseudoscalar field $\phi$ with the Lagrangian
%
\begin{equation}
\label{eq:alp_lagrangian}
\mathcal{L}_\mathrm{ALP}
= -\tfrac{1}{2}(\partial_\mu\phi)^2
  - m_a^2 f_a^2\!\left[1 - \cos\!\left(\frac{\phi}{f_a}\right)\right]
  - \frac{g_{a\gamma}}{4}\,\phi\,F_{\mu\nu}\tilde{F}^{\mu\nu} \,,
\end{equation}
%
where the photon coupling is
%
\begin{equation}
\label{eq:coupling}
g_{a\gamma} = \frac{C_{a\gamma}\,\alpha_\mathrm{em}}{2\pi\,f_a} \,.
\end{equation}
%
We fix two parameters to their natural ECH-motivated values:
%
\begin{itemize}
\item $f_a = M_\mathrm{Pl} = 2.435 \times 10^{18}\;\mathrm{GeV}$
  (the natural scale of the parity-odd sector);
\item $C_{a\gamma} = 8$ (the Standard Model anomaly coefficient from all
  charged fermions: $C = 2\sum_f N_c^{(f)} Q_f^2 = 8$).
\end{itemize}
%
This leaves a two-parameter model with free parameters $\theta_i$ (initial
misalignment angle) and $m_a$ (ALP mass).  The field satisfies
%
\begin{equation}
\ddot{\theta} + 3H\dot{\theta} + m_a^2\sin\theta = 0 \,,
\end{equation}
%
where $\theta \equiv \phi/f_a$ and $H(t)$ is the Hubble parameter in a
standard $\Lambda$CDM background.

\medskip
\noindent\textbf{Spectator regime.}  We require $m_a > \mathrm{few}\times
H_0$, so that the ALP has oscillated and its energy density has redshifted as
$a^{-3}$ by today.  In this regime the ALP contributes negligible energy
density ($\Omega_a \ll 0.01$) and does not affect the background expansion.
Dark energy is provided by a separate cosmological constant~$\Lambda$.  The
spectator regime avoids the rolling-versus-freezing tension identified in
Ref.~\cite{Choi2021}: a single ALP cannot simultaneously provide $w \approx
-1$ dark energy (requiring frozen field) and large birefringence (requiring
rolled field); on the $\Omega_a = 0.68$ contour the maximum achievable
$\beta \approx 0.16^\circ$, a factor of two below observation.


% ---------------------------------------------------------------------------
\subsection{Birefringence Prediction}
\label{sec:beta_prediction}

As the ALP evolves from its initial value $\theta_i$ at recombination to its
present value $\theta_0$, CMB photons accumulate a polarization rotation
%
\begin{equation}
\label{eq:beta}
\beta = \frac{C_{a\gamma}\,\alpha_\mathrm{em}}{4\pi}
        \bigl[\theta(z_\mathrm{rec}) - \theta(z{=}0)\bigr]
      = \frac{C_{a\gamma}\,\alpha_\mathrm{em}\,\theta_i\,\eta}{4\pi} \,,
\end{equation}
%
where the rolling efficiency
$\eta \equiv [\theta(z_\mathrm{rec}) - \theta(0)] / \theta_i$ is computed
numerically by integrating the field equation from $z_\mathrm{rec} = 1089.8$
to $z = 0$ using a DOP853 Runge--Kutta integrator in the variable
$\ln a$~\cite{DOP853}.

A key feature of Eq.~\eqref{eq:beta} is that $f_a$ cancels exactly: $\beta$
depends only on $C_{a\gamma}$, $\alpha_\mathrm{em}$, $\theta_i$, and
$\eta(m_a/H_0)$.  The prediction is therefore robust against
uncertainties in $f_a$ and independent of the specific UV completion.

For the spectator regime ($m_a \gg H_0$, $\eta \to 1$) with $C = 8$ and
$\theta_i = 1$:
%
\begin{equation}
\label{eq:beta_fiducial}
\beta = \frac{8 \times \alpha_\mathrm{em}}{4\pi} \times 1
      \approx 0.27^\circ \,,
\end{equation}
%
in good agreement with the combined measurement $\beta_\mathrm{obs} = 0.342
\pm 0.094^\circ$ from Planck PR4 and ACT
DR6~\cite{Eskilt2022,DiegoPalazuelos2022}.
Figure~\ref{fig:beta_vs_theta} shows the prediction as a function of
$\theta_i$, with the observational band overlaid.


% ---------------------------------------------------------------------------
\subsection{MCMC Constraints}
\label{sec:mcmc}

We constrain the two free parameters $(\theta_i,\,\log_{10}(m_a/\mathrm{eV}))$
using a Gaussian likelihood on the isotropic birefringence angle:
%
\begin{equation}
\ln\mathcal{L} = -\frac{1}{2}
  \left(\frac{\beta_\mathrm{pred} - \beta_\mathrm{obs}}{\sigma_\beta}\right)^2 \,,
\end{equation}
%
with $\beta_\mathrm{obs} = 0.342^\circ$ and $\sigma_\beta = 0.094^\circ$.
We sample with \textsc{Cobaya}~\cite{Cobaya} using 4 MCMC chains, convergence
criterion $\hat{R} - 1 < 0.01$, and flat priors $\theta_i \in [0.01,\,\pi]$,
$\log_{10}(m_a/\mathrm{eV}) \in [-33,\,-28]$.

We note that this likelihood uses the summary statistic $\beta_\mathrm{obs}
\pm \sigma_\beta$ rather than the full $C_\ell^{EB}$ power spectrum employed
in Refs.~\cite{Eskilt2023,Namikawa2025}.  Our analysis therefore does not
constrain the ALP mass through the $\ell$-dependent shape of the EB spectrum
and is less informative for mass discrimination.  However, for the spectator
regime ($m_a \gg H_0$, $\eta \to 1$) the birefringence is
$\ell$-independent, and the summary statistic captures the full information
content.

\medskip
\noindent\textbf{Results.}
After 2160 accepted samples with $\hat{R} - 1 = 0.008$ and effective sample
sizes $N_\mathrm{eff} > 1000$ for both parameters, we obtain:

\begin{table}[h]
\centering
\caption{Posterior summary for the spectator ALP model ($f_a = M_\mathrm{Pl}$,
  $C_{a\gamma} = 8$).}
\label{tab:mcmc}
\begin{tabular}{lcccc}
\hline\hline
Parameter & Mean & Median & 68\% CI & 95\% CI \\
\hline
$\theta_i$ & 1.36 & 1.30 & $[0.90,\, 1.72]$ & $[0.47,\, 2.28]$ \\
$\log_{10}(m_a/\mathrm{eV})$ & $-31.3$ & $-31.2$ & $[-32.4,\, -30.1]$ & $[-32.4,\, -30.0]$ \\
$\beta\;[\mathrm{deg}]$ & 0.336 & 0.335 & $[0.245,\, 0.442]$ & $[0.130,\, 0.519]$ \\
$\eta$ & 0.957 & 0.999 & $[0.95,\, 1.04]$ & $[0.40,\, 1.17]$ \\
\hline\hline
\end{tabular}
\end{table}

\noindent The inferred misalignment angle $\theta_i = 1.36 \pm 0.44$ is
$\mathcal{O}(1)$, consistent with random initial conditions drawn from the
field space $[0,\,2\pi)$.  No fine-tuning is required.  The posterior strongly
favors $\eta \approx 1$ (full rolling), confirming the spectator regime.  The
mass is weakly constrained from below ($m_a \gtrsim 3\,H_0$, where $\eta$
begins to decrease) but unconstrained from above (for $m_a \gg H_0$, $\eta$
saturates at unity and $\beta$ becomes mass-independent).


% ---------------------------------------------------------------------------
\subsection{Model Comparison}
\label{sec:model_comparison}

To assess whether the ALP model adds explanatory value beyond a
phenomenological birefringence parameter, we compare three models:

\begin{enumerate}
\item \textbf{Model~0} ($\beta_\mathrm{free}$): one free parameter $\beta$,
  no physical model.

\item \textbf{Model~2} (ALP, $C = 8$): two free parameters
  $(\theta_i,\,\log_{10} m_a)$; $f_a = M_\mathrm{Pl}$, $C = 8$ fixed.

\item \textbf{Model~2b} (ALP, $C$ free): three free parameters
  $(\theta_i,\,\log_{10} m_a,\,C_{a\gamma})$.
\end{enumerate}

\begin{table}[h]
\centering
\caption{Model comparison.  All non-null models achieve $\chi^2_\mathrm{min}
  \approx 0$ (exact fit to one data point).  $\Delta\chi^2$ is relative to the
  null hypothesis $\beta = 0$.}
\label{tab:model_comparison}
\begin{tabular}{lccccc}
\hline\hline
Model & $k$ & $\chi^2_\mathrm{min}$ & $\Delta\chi^2$ vs null & AIC & $\Delta$AIC vs Model~0 \\
\hline
Null ($\beta = 0$)    & 0 & 13.2 & ---  & ---  & --- \\
Model~0 ($\beta$ free) & 1 & $\approx 0$ & 13.2 (3.6$\sigma$) & 2.0  & 0.0 \\
Model~2 (ALP, $C{=}8$) & 2 & $\approx 0$ & 13.2 (3.6$\sigma$) & 4.0  & $+2.0$ \\
Model~2b (ALP, $C$ free) & 3 & $\approx 0$ & 13.2 (3.6$\sigma$) & 6.0  & $+4.0$ \\
\hline\hline
\end{tabular}
\end{table}

\noindent All three models produce identical $\beta$ posteriors
(Fig.~\ref{fig:beta_comparison}) and identical improvement over the null
hypothesis: $\Delta\chi^2 = 13.2$, corresponding to $3.6\sigma$.  The ALP
model does not fit the data better than free $\beta$.  Information criteria
impose a mild penalty for the additional parameters ($\Delta\mathrm{AIC} =
+2.0$ for Model~2, $+4.0$ for Model~2b).

Floating $C_{a\gamma}$ in Model~2b reveals a degeneracy: the data constrain
the product $C_{a\gamma} \times \theta_i \approx 10.6$ (correlation
$r = -0.52$) but not the individual values.  The posterior width on $C_{a\gamma}$
is 56\% of the prior range, indicating weak constraint.  The Standard Model
value $C = 8$ lies comfortably within the posterior, confirming that no
beyond-SM charged sector is required.


% ---------------------------------------------------------------------------
\subsection{Assessment: Physical Interpretation, Not Statistical Improvement}
\label{sec:honest_assessment}

The model comparison establishes that the spectator ALP provides no
statistical improvement over a phenomenological $\beta$ parameter.  This is
expected: with a single Gaussian measurement, any model that can reproduce
$\beta_\mathrm{obs}$ achieves $\chi^2 \approx 0$; the data constrain exactly
one effective combination regardless of parameterization.

The ALP model's value is physical, not statistical:

\begin{enumerate}
\item \emph{Natural parameters.}  The inferred $\theta_i \approx 1.3$ is
  $\mathcal{O}(1)$, consistent with a randomly drawn initial condition.  A
  phenomenological $\beta = 0.34^\circ$ has no such interpretation.

\item \emph{Structural prediction.}  The model relates $\beta$ to fundamental
  parameters through Eq.~\eqref{eq:beta}: $\beta$ is determined by $C$
  (particle content), $\alpha_\mathrm{em}$ (measured), $\theta_i$ (natural),
  and $\eta$ (computed).  Free $\beta$ has no such structure.

\item \emph{Falsifiability.}  The ALP predicts achromatic, isotropic
  birefringence with a specific mass-dependent evolution.
  LiteBIRD~\cite{LiteBIRD}, with projected $\sigma(\beta) \approx 0.01^\circ$,
  will determine $\theta_i$ to $\pm 0.04\;\mathrm{rad}$, providing a decisive
  test.  If $\beta_\mathrm{pred} = 0.27^\circ$, detection exceeds
  $20\sigma$.  If $\beta = 0$, the model is excluded at comparable
  significance.  The model is also falsified by frequency-dependent
  birefringence or an anisotropic pattern inconsistent with a homogeneous
  scalar background.

\item \emph{Framework context.}  Unlike previous ALP birefringence
  studies~\cite{Fujita2021,Nakagawa2025,LinYanagida2023}, this prediction
  emerges as the sole surviving testable output of a comprehensive
  framework assessment (Sec.~\ref{sec:closure}).  It is not selected from a
  menu of claims but identified by systematic elimination.
\end{enumerate}

We emphasize that spectator ALP models producing similar conclusions have been
studied previously~\cite{Fujita2021,Nakagawa2025}.  Our contribution is the
specific ECH-motivated two-parameter reduction ($f_a = M_\mathrm{Pl}$,
$C = 8$), the explicit comparison against a free-$\beta$ baseline (which, to
our knowledge, has not been performed in the literature), and the framing
within the structural closure of Sec.~\ref{sec:closure}.


% ---------------------------------------------------------------------------
\subsection{Experimental Constraints}
\label{sec:constraints}

Table~\ref{tab:bounds} summarizes existing bounds on the ALP parameter space.
All are satisfied with margins of nine to thirty-two orders of magnitude.

\begin{table}[h]
\centering
\caption{Experimental and observational bounds on the spectator ALP.  The
  model coupling $g_{a\gamma} = 3.8 \times 10^{-21}\;\mathrm{GeV}^{-1}$ and
  mass $m_a \sim 10^{-31}\;\mathrm{eV}$.}
\label{tab:bounds}
\begin{tabular}{lcc}
\hline\hline
Constraint & Bound & Margin \\
\hline
CAST helioscope & $g < 6.6 \times 10^{-11}\;\mathrm{GeV}^{-1}$ & $10^{10}$ \\
SN1987A & $g < 5 \times 10^{-12}\;\mathrm{GeV}^{-1}$ & $10^{9}$ \\
BH superradiance & excludes $6\times10^{-13}\text{--}2\times10^{-11}\;\mathrm{eV}$ & model below \\
CMB spectral distortion & no constraint for $m < 10^{-28}\;\mathrm{eV}$ & model below \\
Isocurvature (Planck) & $f_a < 10^{19}\;\mathrm{GeV}$ for $m \sim H_0$ & marginal \\
DM overproduction & $\theta_i < 2.4$ for $f_a = M_\mathrm{Pl}$ & $\theta_i = 1.3$ safe \\
\hline\hline
\end{tabular}
\end{table}
